\section{Teorijska osnova}
\label{sec:teorija}

\subsection{Korelacija digitalne slike}

Metoda korelacije digitalne slike temelji se na praćenju pomaka
nasumičnog kontrastnog uzorka (engl. speckle pattern),
koji se prije mjerenja nanosi na površinu uzorka, kroz niz slika
snimljenih tijekom deformacije.
Ako je $f(\mathbf{x})$
intenzitet referentne slike u pikselu $\mathbf{x}=(x,y)$,
a $g(\mathbf{x})$ razina
deformirane slike, uz pretpostavku očuvanja intenziteta,
materijalna točka koja se u
referentnom stanju nalazi u $\mathbf{x}$ pomiče se u deformiranom stanju u
$\mathbf{x}+\mathbf{u}(\mathbf{x})$. Tada vrijedi:
\begin{equation}
  f(\mathbf{x}) = g\big(\mathbf{x}+\mathbf{u}(\mathbf{x})\big).
  \label{eq:opt-flow}
\end{equation}
Polje pomaka $\mathbf{u}$ određuje se minimizacijom kvadratnog reziduala intenziteta piksela:
\begin{equation}
  \boldsymbol{\Phi}^2 = \sum_{\mathbf{x}\in\text{ROI}}
  \big[ f(\mathbf{x}) - g(\mathbf{x}+\mathbf{u}(\mathbf{x})) \big]^2 .
  \label{eq:dic-functional}
\end{equation}
U globalnom pristupu korelacije, korištenom u ovom radu, polje pomaka
diskretizira se mrežom konačnih elemenata, dakle pomak unutar elementa interpolira se
funkcijama oblika $N_i$ iz čvornih vrijednosti $\mathbf{u}_i$:
$\mathbf{u}(\mathbf{x})=\sum_i N_i(\mathbf{x})\,\mathbf{u}_i$. Minimizacija
\eqref{eq:dic-functional} tada vodi na linearizirani (Gauss-Newtonov) sustav koji se
rješava iterativno \cite{sutton2009}\cite{hild2006}.

\subsection{Model kamere i projekcijska matrica}

Projekcija točke prostora na ravninu senzora opisuje se tzv. \textit{pinhole} modelom.
Točka $\mathbf{X}=(X,Y,Z)$ u globalnom koordinatnom sustavu
preslikava se u piksel $\mathbf{x}=(x,y)$ preko projekcijske matrice
$\mathsf{P}\in\mathbb{R}^{3\times4}$ u homogenim koordinatama:
\begin{equation}
  s\begin{bmatrix} x \\ y \\ 1 \end{bmatrix}
  = \mathsf{P}\begin{bmatrix} X \\ Y \\ Z \\ 1 \end{bmatrix},
  \qquad
  \mathsf{P} = \mathsf{K}\,[\,\mathsf{R}\ |\ \mathbf{t}\,],
  \label{eq:projection}
\end{equation}
gdje je $s$ skalirni faktor, a $[\mathsf{R}\,|\,\mathbf{t}]$ označava matricu vanjskih (ekstriničnih) parametara kamere koja opisuje rotaciju $\mathsf{R}$ i translaciju $\mathbf{t}$ iz globalnog u koordinatni sustav kamere.

Matrica $\mathsf{K}\in\mathbb{R}^{3\times3}$ predstavlja matricu unutarnjih (intrinzičnih)
parametara kamere, definira se kao:
\begin{equation}
  \mathsf{K} = \begin{bmatrix} f_x & \alpha & c_x \\ 0 & f_y & c_y \\ 0 & 0 & 1 \end{bmatrix},
  \label{eq:intrinsic-matrix}
\end{equation}
gdje su:
\begin{itemize}
  \item $f_x$ i $f_y$ žarišne duljine kamere izražene u pikselima duž horizontalne ($x$) i vertikalne ($y$) osi slikovne ravnine,
  \item $c_x$ i $c_y$ koordinate glavne točke (engl. \textit{principal point}), koja predstavlja sjecište optičke osi kamere i ravnine senzora,
  \item $\alpha$ koeficijent iskošenosti (engl. \textit{skew coefficient}), koji opisuje odstupanje od okomitosti između osi piksela na senzoru (za moderne senzore ova se vrijednost uzima kao nula).
\end{itemize}
Kalibracijom stereovizijskog sustava određuju se projekcijske matrice obiju kamera, $\mathsf{P}_1$ i $\mathsf{P}_2$.

\subsection{Stereo triangulacija}

Ako je poznat piksel $\mathbf{x}_1$ u prvoj i pripadni piksel $\mathbf{x}_2$ u drugoj
kameri, trodimenzijska točka $\mathbf{X}$ rekonstruira se metodom triangulacije. Iz
\eqref{eq:projection} za svaku kameru vrijedi $\mathbf{x}_c \times (\mathsf{P}_c
\tilde{\mathbf{X}})=\mathbf{0}$, što daje po dvije linearne jednadžbe. Slaganjem
jednadžbi obiju kamera dobiva se homogeni sustav
\begin{equation}
  \mathsf{A}\,\tilde{\mathbf{X}} = \mathbf{0}, \qquad
  \mathsf{A} =
  \begin{bmatrix}
    x_1\,\mathbf{p}_1^{3} - \mathbf{p}_1^{1} \\
    y_1\,\mathbf{p}_1^{3} - \mathbf{p}_1^{2} \\
    x_2\,\mathbf{p}_2^{3} - \mathbf{p}_2^{1} \\
    y_2\,\mathbf{p}_2^{3} - \mathbf{p}_2^{2}
  \end{bmatrix},
  \label{eq:dlt}
\end{equation}
gdje je $\mathbf{p}_c^{k}$ $k$-ti redak matrice $\mathsf{P}_c$. Rješenje
$\tilde{\mathbf{X}}$ je desni singularni vektor matrice $\mathsf{A}$ koji odgovara
najmanjoj singularnoj vrijednosti (metoda \textit{Direct Linear Transformation}, DLT)
\cite{hartley2004}.

\subsection{Deformacije na površini}

Iz trodimenzijskog polja pomaka određuju se deformacije. Za svaki trokutasti element
konstruira se lokalni ortonormirani tangencijalni sustav, u kojem se računa gradijent
deformiranja $\mathsf{F}$ te Green-Lagrangeov tenzor deformacije
\begin{equation}
  \mathsf{E} = \tfrac{1}{2}\big(\mathsf{F}^{\mathsf{T}}\mathsf{F} - \mathsf{I}\big).
  \label{eq:green-lagrange}
\end{equation}
Iz komponenti $\mathsf{E}$ dobivaju se glavne deformacije $E_{1,2}$ te ekvivalentna
(von Misesova) deformacija, koje opisuju lokalno deformacijsko stanje površine.
